\section{Permutation rigidity for bases of normal tensors}

When two bases of normal tensors generate the same MPV subspace at every positive length,
the blocks must agree up to a permutation and per-block gauge phases. The
argument turns the equality of spans into an invertible change-of-basis matrix,
then analyses its asymptotics through the overlaps.

\begin{lemma}[Invertible change of basis]
    \label{lem:bnt_invertible_change_basis}
    \lean{exists_invertible_changeBasis}
    \leanok
    Let $v_1, \ldots, v_g$ and $w_1, \ldots, w_g$ be two linearly independent
    families in a complex vector space with the same span. Then there exists an
    invertible matrix $U \in \GL_g(\C)$ such that
    \[
        w_j = \sum_{i=1}^g U_{ij} v_i
    \]
    for every $j$.
\end{lemma}

\begin{proof}\leanok
    Because the two families span the same subspace, each $w_j$ has a unique
    expansion in the basis $\{v_i\}_{i=1}^g$; these coefficients form the
    columns of $U$. If $U a = 0$, then the corresponding linear relation among
    the $w_j$ is trivial because the $w_j$ are linearly independent. Thus $U$
    is injective, hence invertible.
\end{proof}

\begin{lemma}[Eventual change of basis from equal spans]
    \label{lem:bnt_eventual_change_basis}
    \lean{MPSTensor.eventually_exists_invertible_changeBasis}
    \leanok
    \uses{lem:bnt_invertible_change_basis, def:mpv_overlap}
    If two finite families of blocks have asymptotically orthonormal overlaps
    within each family and their MPV spans agree for every system size $N$,
    then for all sufficiently large $N$ there exists an invertible matrix
    $U_N$ expressing the MPV states of one family as linear combinations of
    the MPV states of the other.
\end{lemma}

\begin{proof}\leanok\uses{lem:bnt_overlap_orthonormal_li, lem:bnt_invertible_change_basis}
    By Lemma~\ref{lem:bnt_overlap_orthonormal_li}, both families are
    linearly independent for all sufficiently large $N$. At such an $N$, the
    span equality lets us apply Lemma~\ref{lem:bnt_invertible_change_basis}.
\end{proof}

The change-of-basis matrix is now used to match the blocks. We first treat
equal block counts, then remove that restriction.

\begin{theorem}[Permutation rigidity for equal block counts]
    \label{thm:bnt_perm_same_card}
    \lean{MPSTensor.exists_perm_dimEq_gaugePhaseEquiv_of_overlapOrtho}
    \leanok
    \uses{def:injective, def:mpv_overlap, def:gauge_phase_equiv}
    Let $(A_j)_{j=1}^g$ and $(B_j)_{j=1}^g$ be injective families of positive
    bond dimension satisfying trace-preserving normalization, with self-overlaps tending to~$1$ and
    off-diagonal overlaps tending to~$0$ within each family. If the spans of
    the MPV states agree for every system size $N$, then there is a permutation
    $\pi$ of $\{1, \ldots, g\}$ such that $A_j$ and $B_{\pi(j)}$ have the
    same bond dimension and are gauge-phase equivalent.
\end{theorem}

\begin{proof}\leanok\uses{lem:bnt_eventual_change_basis, thm:overlap_decay, thm:overlap_decay_rect}
    For large $N$, Lemma~\ref{lem:bnt_eventual_change_basis} gives an
    invertible matrix $U_N$ relating the two MPV families. Fix $j$ and write
    $\ket{V^{(N)}(B_j)}=\sum_i u_{ij}^{(N)}\ket{V^{(N)}(A_i)}$. With
    $G_A(N)=(O_{A_iA_\ell}(N))_{i,\ell}$ and
    $b_j(N)=(O_{A_iB_j}(N))_i$, this gives
    \[
        G_A(N)\overline{u_j(N)}=b_j(N),
        \qquad
        O_{B_jB_j}(N)=u_j(N)^{\top}G_A(N)\overline{u_j(N)}.
    \]
    If all mixed overlaps $O_{A_iB_j}(N)$ tended to $0$, then
    $b_j(N)\to0$; since $G_A(N)\to I$, one obtains
    $\overline{u_j(N)}\to0$ and hence $u_j(N)\to0$. Thus
    $O_{B_jB_j}(N)\to0$, contradicting $O_{B_jB_j}(N)\to1$.
    Thus each $B_j$ has a matched block $A_{f(j)}$ with non-decaying mixed
    overlap. Theorems~\ref{thm:overlap_decay_rect}
    and~\ref{thm:overlap_decay} then force equal bond dimension and
    gauge-phase equivalence for each match. Off-diagonal decay within the
    $B$-family makes $f$ injective, hence bijective, and its inverse is the
    required permutation.
\end{proof}

\begin{lemma}[Permutation rigidity with asymptotically orthonormal overlaps]\label{lem:bnt_perm_simple}
    \lean{MPSTensor.exists_eq_numBlocks_and_equiv_gaugePhase_of_overlapOrtho}
    \leanok
    \uses{def:injective, def:mpv, def:mpv_overlap, def:gauge_phase_equiv}
    Let $(A_j)_{j=1}^{g_A}$ and $(B_k)_{k=1}^{g_B}$ be two
    families of injective tensors of positive bond dimension satisfying the TP normalization
    $\sum_i (A_j^i)^\dagger A_j^i = \Id$ and
    $\sum_i (B_k^i)^\dagger B_k^i = \Id$, whose self-overlaps
    converge to~$1$ and whose cross-overlaps within each
    family converge to~$0$.
    If for every system size~$N$,
    \[
        \spn_{\C}\{\ket{V^{(N)}(A_j)} : 1 \le j \le g_A\}
        =
        \spn_{\C}\{\ket{V^{(N)}(B_k)} : 1 \le k \le g_B\},
    \]
    then $g_A = g_B$, and there is a permutation $\pi$ of the common
    block index set such that $A_j$ and $B_{\pi(j)}$ have the same bond
    dimension and are gauge-phase equivalent for each~$j$.
\end{lemma}

\begin{proof}\leanok\uses{lem:bnt_overlap_orthonormal_li, thm:bnt_perm_same_card}
    By Lemma~\ref{lem:bnt_overlap_orthonormal_li}, both families are
    linearly independent for all sufficiently large $N$. At such an $N$, the
    common span has the same dimension when computed from the $A$-family or
    the $B$-family, so $g_A = g_B$. After identifying the two index sets,
    Theorem~\ref{thm:bnt_perm_same_card} yields the permutation conclusion.
\end{proof}

The two preceding results obtain the block matching from the asymptotic
orthonormality of the overlaps. The same matching is obtained by the exact
sector comparison of Section~\ref{sec:sector_bnt_strong_match}, which pairs the
sectors of two canonical forms through a bijection fixed by their bond
dimensions and gauge-phase equivalences. The proof of the Fundamental Theorem
in Chapter~\ref{ch:ft_proof} proceeds through the sector comparison.

The same overlap dichotomy also shows that the MPV states of distinct normal
blocks are non-proportional at arbitrarily large lengths; this is an input to
the exact sector matching later in the chapter.

\begin{theorem}[Non-equivalent irreducible TP blocks are eventually non-proportional]
    \label{thm:not_gauge_phase_eventually_not_proportional}
    \lean{MPSTensor.exists_ge_not_forall_mpv_eq_mul_of_not_gaugePhaseEquiv_of_irreducible_TP}
    \lean{MPSTensor.exists_ge_not_forall_mpv_eq_mul_of_blocksNotGaugePhaseEquiv_of_irreducible_TP}
    \leanok
    \uses{def:gauge_phase_equiv, def:mpv_overlap, def:irreducible_tensor}
    Let $A$ and $B$ be irreducible trace-preserving blocks whose
    self-overlaps converge to $1$. If they are not gauge-phase equivalent,
    then for every lower bound $N_0$ there is $N \ge N_0$ such that
    $\ket{V^{(N)}(A)}$ is not proportional to $\ket{V^{(N)}(B)}$.
    Consequently, in an already separated same-dimensional BNT block family,
    any two distinct blocks have such non-proportional MPV states at
    arbitrarily large lengths.
\end{theorem}

\begin{proof}\leanok
    Suppose that for all sufficiently large $N$ one had
    $\ket{V^{(N)}(A)}=c_N\ket{V^{(N)}(B)}$. Then
    \begin{equation}\label{eq:bnt_proportional_overlap_identities}
        O_{AA}(N)=|c_N|^2O_{BB}(N),
        \qquad
        O_{AB}(N)=c_NO_{BB}(N).
    \end{equation}
    Since $O_{AA}(N)\to1$ and $O_{BB}(N)\to1$, the first identity in
    (\ref{eq:bnt_proportional_overlap_identities}) gives $|c_N|\to1$, and hence
    $|O_{AB}(N)|\to1$. The irreducible
    trace-preserving overlap decay theorem gives $O_{AB}(N)\to0$ when
    $A\not\sim_{\mathrm{gp}}B$, a contradiction.
\end{proof}

\section{Newton--Girard identities and power-sum recovery}

The coefficients $c_N^{(j)} = \sum_q \mu_{j,q}^N$ that appear in the
two-layer BNT expansion are power sums of the copy weights. Recovering the
individual weights from these sums is a Newton--Girard problem: equal power sums
determine equal characteristic polynomials, hence equal multisets of roots.

\begin{theorem}[Newton--Girard trace recursion]\label{thm:newton_girard}
    \lean{Matrix.charpoly_eq_of_forall_trace_pow_eq,
        Matrix.charpoly_eq_of_trace_pow_eq_of_le_card}
    \leanok
    If $A, B \in \MN{n}$ satisfy
    $\tr(A^k) = \tr(B^k)$ for all $k \ge 1$,
    then $A$ and $B$ have the same characteristic polynomial.
    It is enough to assume this equality for $1 \le k \le n$.
\end{theorem}

\begin{proof}
    \leanok
    Let $e_m$ denote the $m$th elementary symmetric polynomial in the
    eigenvalues, with $e_0 = 1$.
    The Newton--Girard recursion
    \[
        m e_m = \sum_{i=1}^m (-1)^{i-1} e_{m-i} \tr(A^i)
    \]
    expresses the coefficients of the characteristic polynomial in terms
    of the power-sum traces $\tr(A^i)$.
    Equal power sums for $A$ and $B$ therefore give equal coefficients,
    hence equal characteristic polynomials.
\end{proof}

The matched-sector comparison below applies the power-sum identities to a single
matched pair, after the BNT permutation and phase gauges have identified which
basis tensors are compared: the coefficient $\sum_q \mu_{j,q}^N$ then equals
$\sum_q (\nu_{j,q} e^{i\phi_j})^N$, and Newton--Girard recovers the weight
multiset with multiplicity.

\begin{lemma}[Eventual vanishing of a finite geometric sum]
    \label{lem:finite_geometric_sum_eventual_vanishing}
    \lean{MPSTensor.SectorWeightData.geom_sum_eventually_zero}
    \leanok
    Let $w_1,\ldots,w_n$ be nonzero complex numbers.  If
    \[
        \sum_i c_i w_i^N=0
    \]
    for every sufficiently large $N$, then the same equality holds for every
    nonnegative exponent.
\end{lemma}

\begin{proof}\leanok
    Induct on the number of terms.  Subtracting $w_1$ times the equality at
    exponent $N$ from the equality at exponent $N+1$ gives
    \[
        \sum_{i=2}^{n} c_i(w_i-w_1)w_i^N=0.
    \]
    The induction hypothesis yields this equality at every exponent, so the
    original sum $S_N=\sum_i c_iw_i^N$ satisfies $S_{N+1}=w_1S_N$ for every
    $N$.  Since $S_N$ eventually vanishes and $w_1\ne0$, it follows that
    $S_0=0$ and hence $S_N=0$ for every $N$.
\end{proof}

\begin{lemma}[Eventual power-sum equality holds at every exponent]
    \label{lem:eventual_power_sum_equality_all_exponents}
    \lean{MPSTensor.SectorWeightData.power_sums_eq_of_eventually_eq_hetero}
    \leanok
    \uses{lem:finite_geometric_sum_eventual_vanishing}
    Let $a_1,\ldots,a_m$ and $b_1,\ldots,b_n$ be nonzero complex numbers.  If
    their power sums agree for every sufficiently large exponent, then
    \[
        \sum_i a_i^k=\sum_j b_j^k
    \]
    for every nonnegative exponent $k$.
\end{lemma}

\begin{proof}\leanok
    Set $w_l=a_l$ for $l\leq m$ and $w_{m+l}=b_l$ for $l\leq n$, with
    coefficients $c_l=1$ for $l\leq m$ and $c_l=-1$ for $l>m$.  Then
    \[
        \sum_{l=1}^{m+n}c_lw_l^k
        =\sum_i a_i^k-\sum_j b_j^k.
    \]
    The hypothesis makes this expression zero for every sufficiently large
    $k$.  Lemma~\ref{lem:finite_geometric_sum_eventual_vanishing} makes it zero
    for every $k$, which is the required power-sum equality.
\end{proof}

\begin{theorem}[Finite-range equal power sums imply equal multisets]
    \label{thm:bounded_power_sum_multiset}
    \lean{Matrix.sum_pow_eq_implies_multiset_eq_of_le_card,
        Matrix.sum_pow_eq_implies_nonzero_multiset_eq_of_le_max_length,
        Matrix.sum_pow_eq_implies_multiset_eq_of_le_max_length,
        Matrix.sum_pow_eq_implies_card_eq_and_multiset_eq_of_le_max_card}
    \leanok
    Let $\alpha, \beta : \{0, \ldots, n-1\} \to \C$. If
    $\sum_i \alpha_i^k = \sum_i \beta_i^k$ for $1 \le k \le n$, then the
    multisets $\{\alpha_i\}$ and $\{\beta_i\}$ are equal.

    Two variants relax the equal-cardinality assumption. If
    $\alpha : \{0, \ldots, m-1\} \to \C$ and $\beta : \{0, \ldots, n-1\} \to \C$
    have no zero entries and their power sums agree for $1 \le k \le \max(m,n)$,
    then $m=n$ and the multisets are equal. Without a nonzero hypothesis, the
    list form $\sum_{k=1}^{x_a} \lambda_{a,k}^N = \sum_{k=1}^{x_b} \lambda_{b,k}^N$
    for $1 \le N \le \max(x_a,x_b)$ gives equality of the nonzero multisets; if
    it also holds at $N=0$ then $x_a=x_b$ and the full lists agree.
\end{theorem}

\begin{proof}
    \leanok
    \uses{thm:newton_girard}
    Apply the finite-range Newton--Girard theorem to $\diag(\alpha)$ and
    $\diag(\beta)$ and extract the roots of $\prod_i (X - \alpha_i)$ and
    $\prod_i (X - \beta_i)$. For unequal cardinalities, pad both lists with zeros
    to length $\max(m,n)$; the nonzero hypothesis makes the padded-zero counts
    $\max(m,n)-m$ and $\max(m,n)-n$, forcing $m=n$. Without that hypothesis,
    delete the zero entries first (unchanged for positive $N$); the equality at
    $N=0$ gives $x_a=x_b$.
\end{proof}

\begin{corollary}[A nonzero finite family has a nonvanishing power sum]
    \label{cor:nonvanishing_finite_power_sum}
    \lean{Matrix.exists_sum_pow_ne_zero_of_pos_card}
    \leanok
    \uses{thm:bounded_power_sum_multiset}
    Let $m\geq 1$ and let $\lambda_0,\ldots,\lambda_{m-1}\in\C$ be nonzero.
    Then there is an exponent $N$ with $1\leq N\leq m$ such that
    \[
        \sum_{q=0}^{m-1}\lambda_q^N\neq 0.
    \]
    This is the consequence of the Appendix power-sum lemma used in the proof
    of Lemma~L
    \cite[Appendix~C.3, lines~1851--1858]{Cirac2016MPDO_arXiv}; the finite-range
    comparison is
    \cite[Appendix~A, lines~1155--1163]{Cirac2016MPDO_arXiv}.
\end{corollary}

\begin{proof}
    \leanok
    \uses{thm:bounded_power_sum_multiset}
    Suppose that every power sum with $1\leq N\leq m$ vanishes. Comparing
    $\lambda_0,\ldots,\lambda_{m-1}$ with the empty family in
    Theorem~\ref{thm:bounded_power_sum_multiset} gives $m=0$, contradicting
    $m\geq 1$.
\end{proof}

\begin{theorem}[Equal power sums imply equal multisets]\label{thm:power_sum_multiset}
    \lean{Matrix.sum_pow_eq_implies_multiset_eq}
    \leanok
    Let $\alpha, \beta : \{0, \ldots, n-1\} \to \C$.
    If $\sum_i \alpha_i^k = \sum_i \beta_i^k$ for all $k \ge 1$,
    then the multisets $\{\alpha_i\}$ and $\{\beta_i\}$ are equal.
\end{theorem}

\begin{proof}
    \leanok
    \uses{thm:newton_girard}
    Construct diagonal matrices $\diag(\alpha)$ and
    $\diag(\beta)$.
    The power-sum hypothesis gives
    $\tr(\diag(\alpha)^k) = \tr(\diag(\beta)^k)$.
    By Theorem~\ref{thm:newton_girard}, the characteristic
    polynomials agree: $\prod_i (X - \alpha_i) = \prod_i (X - \beta_i)$.
    Extracting roots gives multiset equality.
\end{proof}

% Maintainer note: Theorems thm:newton_girard and thm:power_sum_multiset feed the
% equal-MPV corollary of Section~\ref{sec:ft_equal_proportional_cases}, where the
% multiplicities r_j and weights \mu_{j,q} are recovered; the BNT multiplicity
% lemmas handle that recovery directly.
